\chapter{Graphs: Basics and Traversals}
\chapterepigraph{A graph is nothing more than a collection of states and
the moves that connect them. Once a problem's ``states'' and ``valid
moves'' are identified, traversing the graph -- visiting every reachable
state exactly once -- answers a surprising number of questions that
don't look like graph problems at first glance.}

\section{Representing a Graph, and the Two Ways to Walk It}
\label{sec:graph-basics-intro}

A graph is a set of \textbf{nodes} (vertices) and \textbf{edges}
(connections between nodes, which may be directed or undirected,
weighted or unweighted). The standard representation for sparse graphs
(the vast majority of interview problems) is an \textbf{adjacency
list}: an array (or hash map) where index $i$ holds a list of every
node directly reachable from node $i$.

\begin{patternbox}[Recognizing a Graph Problem]
A problem is a graph problem in disguise whenever it describes
\textbf{states} and \textbf{transitions between states} -- grid cells
and their neighbors, people and their connections, words and
single-character edits, courses and their prerequisites. The moment a
problem can be rephrased as ``starting from X, what can I reach by
repeatedly taking valid moves,'' a graph traversal applies.
\end{patternbox}

\begin{ideabox}[Depth-First Search vs. Breadth-First Search]
Both traversals visit every node reachable from a start node exactly
once, using a $\textit{visited}$ array to avoid repeats (and infinite
loops in cyclic graphs) -- they differ only in \textbf{which data
structure} drives the order of exploration:
\begin{itemize}
  \item \textbf{DFS} uses a \textbf{stack} (often the recursion call
        stack itself): go as deep as possible down one path before
        backtracking. Natural for problems about \emph{connectivity}
        and \emph{exhaustive exploration} (does a path exist, what is
        everything reachable).
  \item \textbf{BFS} uses an explicit \textbf{queue}: explore all
        nodes at distance $1$ before any node at distance $2$. Natural
        for problems about \emph{shortest distance in an unweighted
        graph} or \emph{level-by-level} processing (multi-source
        spreading, like fire or rot spreading outward).
\end{itemize}
\end{ideabox}

\begin{patternbox}[What This Chapter Covers]
\begin{itemize}
  \item \textbf{DFS Traversal (Section~\ref{sec:graph-dfs})} and
        \textbf{BFS Traversal (Section~\ref{sec:graph-bfs})}: the two
        foundational walks, established once and reused by name in
        every later chapter.
  \item \textbf{Number of Provinces (Section~\ref{sec:graph-number-of-provinces})}:
        counting connected components -- run a traversal from every
        unvisited node, counting how many traversals it takes.
  \item \textbf{Clone a Graph (Section~\ref{sec:graph-clone-graph})}:
        a traversal that builds a parallel copy of the graph as it
        visits, using a hash map to avoid duplicating already-cloned
        nodes.
  \item \textbf{Rotten Oranges (Section~\ref{sec:graph-rotten-oranges})},
        \textbf{Flood Fill (Section~\ref{sec:graph-flood-fill})}, and
        \textbf{01 Matrix (Section~\ref{sec:graph-01-matrix})}: treating
        a grid as an implicit graph (each cell a node, its $4$ neighbors
        its edges), with the first and last introducing
        \textbf{multi-source BFS} -- starting a single BFS from
        \emph{every} source simultaneously.
  \item \textbf{Number of Islands (Section~\ref{sec:graph-number-of-islands})}:
        Number of Provinces' connected-component counting, applied to
        an implicit grid graph instead of an explicit adjacency list.
\end{itemize}
\end{patternbox}

\newpage

\section{Depth-First Search (DFS) Traversal}
\label{sec:graph-dfs}

\problemstatement{Given a graph (as an adjacency list) and a starting
node, visit every node reachable from the start exactly once, exploring
as deep as possible along each path before backtracking.}

\begin{patternbox}
\emph{``Visit every reachable node, going deep before wide''} is solved
by recursion (or an explicit stack): visit the current node, mark it
visited, then recurse into every unvisited neighbor in turn.
\end{patternbox}

\subsection*{Algorithm}

\begin{enumerate}
  \item Maintain a $\textit{visited}$ array, initialized to all false.
  \item Define $\textit{dfs}(\textit{node})$: mark $\textit{node}$ as
        visited, add it to the result; for every neighbor of
        $\textit{node}$ not yet visited, recursively call
        $\textit{dfs}(\textit{neighbor})$.
  \item To traverse the whole graph (which may be disconnected), call
        $\textit{dfs}(i)$ for every node $i$ not yet visited.
\end{enumerate}

\subsection*{Why a \textit{visited} array is required}

Without tracking visited nodes, a cycle in the graph (or even a simple
back-and-forth edge between two nodes) would cause infinite recursion.
Marking a node visited \textbf{before} recursing into its neighbors
(not after returning) is essential: it guarantees that if two different
neighbors both point back to an already-in-progress node, only the
first one to reach it recurses into it.

\subsection*{C++ implementation}

\begin{lstlisting}
#include <bits/stdc++.h>
using namespace std;

void dfs(int node, vector<vector<int>>& adj, vector<bool>& visited,
         vector<int>& result) {
    visited[node] = true;
    result.push_back(node);

    for (int neighbor : adj[node]) {
        if (!visited[neighbor]) {
            dfs(neighbor, adj, visited, result);
        }
    }
}

vector<int> dfsTraversal(int n, vector<vector<int>>& adj) {
    vector<bool> visited(n, false);
    vector<int> result;

    for (int i = 0; i < n; i++) {
        if (!visited[i]) {
            dfs(i, adj, visited, result);
        }
    }
    return result;
}

int main() {
    int n = 5;
    vector<vector<int>> adj(n);
    auto addEdge = [&](int u, int v) { adj[u].push_back(v); adj[v].push_back(u); };
    addEdge(0, 1); addEdge(0, 2); addEdge(1, 3); addEdge(2, 4);

    for (int x : dfsTraversal(n, adj)) cout << x << " ";
    cout << endl; // 0 1 3 2 4
    return 0;
}
\end{lstlisting}

\subsection*{Dry run}

\dryrun{Take the graph with edges $0{-}1$, $0{-}2$, $1{-}3$, $2{-}4$,
starting traversal from node $0$.}

Visit $0$ (mark visited, record). Its neighbors: $1$ (unvisited) --
recurse. Visit $1$, record. Its neighbors: $0$ (visited, skip), $3$
(unvisited) -- recurse. Visit $3$, record; no unvisited neighbors,
return. Back at $1$: no more neighbors, return. Back at $0$: next
neighbor $2$ (unvisited) -- recurse. Visit $2$, record. Its neighbors:
$0$ (visited, skip), $4$ (unvisited) -- recurse. Visit $4$, record; no
unvisited neighbors, return. Final order: $0,1,3,2,4$.

\begin{complexitybox}
\textbf{Time Complexity.} $O(V+E)$ -- every vertex is visited once
(marking it visited), and every edge is examined once (from each
endpoint, checking if the other side is visited).
\textbf{Space Complexity.} $O(V)$ for the \textit{visited} array and
the recursion stack (which can be as deep as $V$ in a path-shaped
graph).
\end{complexitybox}

\begin{takeawaybox}
DFS's ``visit, mark, recurse into unvisited neighbors'' template is the
single most reused piece of code in this chapter -- nearly every later
problem either calls this function directly or modifies one small part
of it (the visiting action, the stopping condition, or what gets
collected).
\end{takeawaybox}

\section{Breadth-First Search (BFS) Traversal}
\label{sec:graph-bfs}

\problemstatement{Given a graph (as an adjacency list) and a starting
node, visit every node reachable from the start exactly once,
exploring all nodes at distance $1$ before any node at distance $2$,
and so on.}

\begin{patternbox}
\emph{``Visit level by level, nearest first''} is solved with a
\textbf{queue}: push the start node, then repeatedly pop the front,
visit it, and push every unvisited neighbor onto the back -- the
first-in-first-out order of a queue naturally processes nodes in
increasing distance from the start.
\end{patternbox}

\subsection*{Algorithm}

\begin{enumerate}
  \item Maintain a $\textit{visited}$ array; mark the start node
        visited and push it onto a queue.
  \item While the queue is non-empty: pop the front node, add it to
        the result; for every unvisited neighbor, mark it visited
        \textbf{immediately} (before it is popped) and push it onto
        the queue.
  \item To traverse the whole graph, repeat for every node not yet
        visited.
\end{enumerate}

\begin{ideabox}[Mark Visited on Push, Not on Pop]
A subtle but important detail: a node must be marked visited the
moment it is \textbf{pushed} onto the queue, not when it is popped.
If marking were delayed until pop time, the same node could be pushed
multiple times by different neighbors before any of those pushes is
processed, wasting work and (in graphs with cycles) potentially
causing incorrect results.
\end{ideabox}

\subsection*{C++ implementation}

\begin{lstlisting}
#include <bits/stdc++.h>
using namespace std;

vector<int> bfsFromSource(int src, vector<vector<int>>& adj, vector<bool>& visited) {
    vector<int> result;
    queue<int> q;
    visited[src] = true;
    q.push(src);

    while (!q.empty()) {
        int node = q.front();
        q.pop();
        result.push_back(node);

        for (int neighbor : adj[node]) {
            if (!visited[neighbor]) {
                visited[neighbor] = true;
                q.push(neighbor);
            }
        }
    }
    return result;
}

vector<int> bfsTraversal(int n, vector<vector<int>>& adj) {
    vector<bool> visited(n, false);
    vector<int> result;

    for (int i = 0; i < n; i++) {
        if (!visited[i]) {
            vector<int> part = bfsFromSource(i, adj, visited);
            result.insert(result.end(), part.begin(), part.end());
        }
    }
    return result;
}

int main() {
    int n = 5;
    vector<vector<int>> adj(n);
    auto addEdge = [&](int u, int v) { adj[u].push_back(v); adj[v].push_back(u); };
    addEdge(0, 1); addEdge(0, 2); addEdge(1, 3); addEdge(2, 4);

    for (int x : bfsTraversal(n, adj)) cout << x << " ";
    cout << endl; // 0 1 2 3 4
    return 0;
}
\end{lstlisting}

\subsection*{Dry run}

\dryrun{Take the same graph as Section~\ref{sec:graph-dfs}: edges
$0{-}1$, $0{-}2$, $1{-}3$, $2{-}4$, starting from node $0$.}

Queue: $[0]$ (visited $=\{0\}$). Pop $0$, record; push unvisited
neighbors $1,2$: queue $[1,2]$ (visited $=\{0,1,2\}$). Pop $1$,
record; push unvisited neighbor $3$: queue $[2,3]$. Pop $2$, record;
push unvisited neighbor $4$: queue $[3,4]$. Pop $3$, record; no new
neighbors. Pop $4$, record; no new neighbors. Final order:
$0,1,2,3,4$ -- notice this differs from DFS's $0,1,3,2,4$: BFS visits
both of $0$'s direct neighbors ($1$ and $2$) before descending to
either of their children, exactly the ``nearest first'' property.

\begin{complexitybox}
\textbf{Time Complexity.} $O(V+E)$ -- identical bound to DFS, for the
same reason (every vertex and edge examined once).
\textbf{Space Complexity.} $O(V)$ for the \textit{visited} array and
the queue (which can hold up to all $V$ nodes in the worst case, e.g.\
a star graph).
\end{complexitybox}

\begin{takeawaybox}
DFS and BFS visit the exact same set of reachable nodes and share the
same $O(V+E)$ complexity -- the only difference is traversal
\textbf{order}, governed entirely by swapping a stack for a queue.
Whenever a problem specifically needs shortest paths in an unweighted
graph or strictly level-by-level processing, that need points directly
at BFS; everything else defaults naturally to whichever is more
convenient to implement (usually DFS, via plain recursion).
\end{takeawaybox}

\section{Number of Provinces}
\label{sec:graph-number-of-provinces}

\problemstatement{Given an $n \times n$ adjacency matrix
$\textit{isConnected}$ for $n$ cities (where
$\textit{isConnected}[i][j]=1$ means cities $i$ and $j$ are directly
connected), find the number of \textbf{provinces} -- groups of cities
connected directly or indirectly, where a city in one group has no
connection to a city in any other group.}

\begin{patternbox}
\emph{``Count the number of disconnected groups''} is exactly counting
\textbf{connected components}: run a traversal (DFS or BFS) from every
\textbf{unvisited} node, marking everything that traversal reaches.
Each time a traversal is launched from a fresh, unvisited node, that is
one new province discovered.
\end{patternbox}

\subsection*{Algorithm}

\begin{enumerate}
  \item Maintain a $\textit{visited}$ array and a $\textit{provinces}$
        counter, both initialized to $0$/false.
  \item For each city $i$ from $0$ to $n-1$: if $i$ is not visited,
        increment $\textit{provinces}$, then run a full DFS (or BFS)
        from $i$, marking every reachable city visited.
  \item Return $\textit{provinces}$.
\end{enumerate}

\begin{ideabox}[Why This Always Terminates with the Correct Count]
Every city belongs to exactly one province (provinces are, by
definition, a partition of the cities). The traversal from an
unvisited city visits \emph{exactly} that city's entire province (no
more, no less, since the adjacency matrix's connections define the
province's boundary). So each outer-loop iteration that finds an
unvisited city is necessarily the \emph{first} city of a province not
yet counted -- which is precisely why incrementing the counter there,
and nowhere else, gives the correct total.
\end{ideabox}

\subsection*{C++ implementation}

\begin{lstlisting}
#include <bits/stdc++.h>
using namespace std;

void dfs(int node, vector<vector<int>>& isConnected, vector<bool>& visited) {
    visited[node] = true;
    int n = isConnected.size();
    for (int neighbor = 0; neighbor < n; neighbor++) {
        if (isConnected[node][neighbor] == 1 && !visited[neighbor]) {
            dfs(neighbor, isConnected, visited);
        }
    }
}

int findCircleNum(vector<vector<int>>& isConnected) {
    int n = isConnected.size();
    vector<bool> visited(n, false);
    int provinces = 0;

    for (int i = 0; i < n; i++) {
        if (!visited[i]) {
            provinces++;
            dfs(i, isConnected, visited);
        }
    }
    return provinces;
}

int main() {
    vector<vector<int>> isConnected = {
        {1, 1, 0},
        {1, 1, 0},
        {0, 0, 1}
    };
    cout << findCircleNum(isConnected) << endl; // 2
    return 0;
}
\end{lstlisting}

\subsection*{Dry run}

\dryrun{Take the $3\times3$ matrix above (cities $0,1$ connected;
city $2$ isolated).}

$i=0$: unvisited, $\textit{provinces}=1$; DFS from $0$ visits $1$
(connected) and stops (city $2$ not connected to $0$ or $1$);
visited $=\{0,1\}$. $i=1$: already visited, skip. $i=2$: unvisited,
$\textit{provinces}=2$; DFS from $2$ visits only itself. Final answer:
$2$ provinces.

\begin{complexitybox}
\textbf{Time Complexity.} $O(n^2)$ -- the adjacency matrix
representation forces an $O(n)$ scan per node during DFS (checking
every possible neighbor), across $n$ nodes total.
\textbf{Space Complexity.} $O(n)$ for the \textit{visited} array and
recursion stack.
\end{complexitybox}

\begin{takeawaybox}
``Count the number of connected components'' is one of the most
frequently recurring graph templates: launch a traversal from every
unvisited node, counting launches. This exact loop structure --
\textit{not} the traversal itself -- is what changes across this
chapter's later problems (Number of Islands, Section~\ref{sec:graph-number-of-islands},
reuses this identical loop on an implicit grid graph instead of an
explicit matrix).
\end{takeawaybox}

\section{Clone a Graph}
\label{sec:graph-clone-graph}

\problemstatement{Given a reference to a node in a connected
undirected graph (where each node stores a value and a list of
neighbor pointers), return a \textbf{deep copy} (clone) of the entire
graph -- a structurally identical graph made of entirely new node
objects.}

\begin{patternbox}
\emph{``Copy every node and every connection, visiting each node
exactly once''} is a traversal (DFS or BFS) augmented with a
\textbf{hash map from original node to its clone} -- the map serves
double duty: it remembers which nodes have already been cloned (acting
as the traversal's \textit{visited} check) and it provides the correct
clone object to link into a neighbor's adjacency list.
\end{patternbox}

\subsection*{Why a hash map, not a simple \textit{visited} array}

Every traversal in this chapter needs to track which nodes have been
visited -- but previous problems used array indices (city numbers,
grid coordinates) as keys, which a simple boolean array handles
directly. Here, nodes are objects with no natural array index, and
\emph{visiting} a node is inseparable from \emph{creating its clone} --
so a hash map from original node pointer to cloned node pointer plays
both roles: checking $\textit{map.find(node)} \neq \textit{map.end()}$
is the visited check, and $\textit{map}[\textit{node}]$ retrieves the
already-built clone to attach as a neighbor.

\subsection*{Algorithm (DFS)}

\begin{enumerate}
  \item Define $\textit{cloneNode}(\textit{node})$: if $\textit{node}$
        is already in the map, return its clone immediately (this
        breaks cycles -- without this check, a cyclic graph would
        recurse forever).
  \item Otherwise, create a new node with the same value, store it in
        the map \textbf{before} recursing into neighbors (exactly like
        marking visited before recursing in plain DFS), then for each
        original neighbor, recursively clone it and append the result
        to the new node's neighbor list.
  \item Return $\textit{cloneNode}(\textit{startNode})$.
\end{enumerate}

\subsection*{C++ implementation}

\begin{lstlisting}
#include <bits/stdc++.h>
using namespace std;

struct Node {
    int val;
    vector<Node*> neighbors;
    Node(int v) : val(v) {}
};

Node* cloneNode(Node* node, unordered_map<Node*, Node*>& cloneMap) {
    if (cloneMap.find(node) != cloneMap.end()) {
        return cloneMap[node];
    }

    Node* clone = new Node(node->val);
    cloneMap[node] = clone;

    for (Node* neighbor : node->neighbors) {
        clone->neighbors.push_back(cloneNode(neighbor, cloneMap));
    }
    return clone;
}

Node* cloneGraph(Node* startNode) {
    if (startNode == nullptr) return nullptr;
    unordered_map<Node*, Node*> cloneMap;
    return cloneNode(startNode, cloneMap);
}

int main() {
    Node* a = new Node(1);
    Node* b = new Node(2);
    Node* c = new Node(3);
    a->neighbors = {b, c};
    b->neighbors = {a, c};
    c->neighbors = {a, b};

    Node* clonedA = cloneGraph(a);
    cout << clonedA->val << " has " << clonedA->neighbors.size()
         << " neighbors" << endl; // 1 has 2 neighbors
    cout << (clonedA == a ? "same object" : "different object") << endl; // different object
    return 0;
}
\end{lstlisting}

\subsection*{Dry run}

\dryrun{Take a triangle graph: nodes $1,2,3$, each connected to the
other two. Start cloning from node $1$.}

$\textit{cloneNode}(1)$: not in map, create clone $1'$, store
$\textit{map}[1]=1'$. Recurse into neighbor $2$:
$\textit{cloneNode}(2)$: not in map, create clone $2'$, store
$\textit{map}[2]=2'$. Recurse into $2$'s neighbors, $1$ and $3$.
$\textit{cloneNode}(1)$: \textbf{already in map}, return $1'$
immediately (no infinite recursion, even though $1$ and $2$ point back
to each other). $\textit{cloneNode}(3)$: not in map, create clone
$3'$, store $\textit{map}[3]=3'$; recurse into $3$'s neighbors $1,2$,
both already in map, return $1',2'$ immediately. Back at $2'$:
neighbors $=[1',3']$. Back at $1'$: continue to neighbor $3$, already
in map, return $3'$. Final: $1'$'s neighbors $=[2',3']$ -- a complete,
independent clone of the original triangle.

\begin{complexitybox}
\textbf{Time Complexity.} $O(V+E)$ -- every node is cloned exactly
once (guarded by the map check), and every edge is examined exactly
once (from each direction).
\textbf{Space Complexity.} $O(V)$ for the hash map and the recursion
stack.
\end{complexitybox}

\begin{takeawaybox}
Whenever a traversal needs to \textbf{build something new} while
visiting (a clone, a transformed copy, a mapping to a different
representation) rather than simply recording or counting, a hash map
keyed by the original node typically replaces the plain
\textit{visited} array -- the map's presence-check still prevents
revisiting and infinite loops, while also handing back the
already-built result the moment it's needed again.
\end{takeawaybox}

\section{Rotten Oranges}
\label{sec:graph-rotten-oranges}

\problemstatement{Given a grid where each cell is $0$ (empty), $1$
(fresh orange), or $2$ (rotten orange), every minute any fresh orange
\textbf{4-directionally adjacent} to a rotten orange becomes rotten.
Find the minimum number of minutes until no fresh orange remains, or
$-1$ if some fresh orange can never be reached.}

\begin{patternbox}
\emph{``Multiple sources spread outward simultaneously, find the time
for the spread to finish''} is solved by \textbf{multi-source BFS}: push
\textbf{every} initially rotten orange into the queue at the start
(not just one), then run ordinary BFS -- the key realization is that
BFS already processes nodes in order of distance from \emph{whichever}
source is nearest, automatically, when all sources start in the queue
together at time $0$.
\end{patternbox}

\subsection*{Why starting from all sources at once is correct}

If BFS were run separately from each rotten orange one at a time, the
correctness would still hold but the bookkeeping (combining several
separate distance computations, taking a minimum over all of them)
would be needlessly complex. Pushing every source into the \emph{same}
queue at the very start works because BFS's queue already guarantees
nodes are processed in non-decreasing distance order \emph{from
whichever source reaches them first} -- exactly the quantity this
problem wants (a fresh orange rots at the minute the \emph{nearest}
rotten orange's spread reaches it).

\subsection*{Algorithm}

\begin{enumerate}
  \item Scan the grid: push every rotten orange's coordinates into a
        queue, and count the total number of fresh oranges.
  \item Run BFS level by level (processing the entire current queue
        before moving to the next level represents one minute
        passing): for each rotten orange popped, check its $4$
        neighbors; any fresh neighbor becomes rotten, decrement the
        fresh count, and push it for the next level.
  \item Track the number of levels (minutes) processed. After BFS
        completes, if any fresh oranges remain uncounted, return $-1$;
        otherwise return the number of minutes elapsed.
\end{enumerate}

\subsection*{C++ implementation}

\begin{lstlisting}
#include <bits/stdc++.h>
using namespace std;

int orangesRotting(vector<vector<int>>& grid) {
    int rows = grid.size(), cols = grid[0].size();
    queue<pair<int,int>> q;
    int freshCount = 0;

    for (int r = 0; r < rows; r++) {
        for (int c = 0; c < cols; c++) {
            if (grid[r][c] == 2) q.push({r, c});
            else if (grid[r][c] == 1) freshCount++;
        }
    }

    int minutes = 0;
    int dr[] = {-1, 1, 0, 0};
    int dc[] = {0, 0, -1, 1};

    while (!q.empty() && freshCount > 0) {
        int levelSize = q.size();
        bool rottedThisLevel = false;

        for (int i = 0; i < levelSize; i++) {
            auto [r, c] = q.front();
            q.pop();

            for (int dir = 0; dir < 4; dir++) {
                int nr = r + dr[dir], nc = c + dc[dir];
                if (nr >= 0 && nr < rows && nc >= 0 && nc < cols &&
                    grid[nr][nc] == 1) {
                    grid[nr][nc] = 2;
                    freshCount--;
                    q.push({nr, nc});
                    rottedThisLevel = true;
                }
            }
        }
        if (rottedThisLevel) minutes++;
    }
    return freshCount == 0 ? minutes : -1;
}

int main() {
    vector<vector<int>> grid = {
        {2, 1, 1},
        {1, 1, 0},
        {0, 1, 1}
    };
    cout << orangesRotting(grid) << endl; // 4
    return 0;
}
\end{lstlisting}

\subsection*{Dry run}

\dryrun{Take the $3\times3$ grid above (one rotten orange at
$(0,0)$).}

Initial queue: $[(0,0)]$, $\textit{freshCount}=6$. Minute $1$: pop
$(0,0)$, rot neighbors $(1,0)$ and $(0,1)$ (both fresh); push both;
$\textit{freshCount}=4$. Minute $2$: pop $(1,0)$, rot $(1,1)$; pop
$(0,1)$, rot $(0,2)$ ($(1,1)$ is already rotten by this point, so
$(0,1)$'s other neighbor is skipped); $\textit{freshCount}=2$. Minute
$3$: pop $(1,1)$, rot $(2,1)$; pop $(0,2)$, no fresh neighbors;
$\textit{freshCount}=1$. Minute $4$: pop $(2,1)$, rot $(2,2)$;
$\textit{freshCount}=0$. The queue still holds $(2,2)$, but the loop
exits immediately since $\textit{freshCount}=0$ -- final answer: $4$
minutes.

\begin{complexitybox}
\textbf{Time Complexity.} $O(\textit{rows} \times \textit{cols})$ --
every cell is pushed onto the queue and processed at most once.
\textbf{Space Complexity.} $O(\textit{rows} \times \textit{cols})$ for
the queue in the worst case (a grid that is entirely rotten oranges
initially).
\end{complexitybox}

\begin{takeawaybox}
Multi-source BFS -- seeding the queue with every source simultaneously
rather than running BFS once per source -- is the standard technique
whenever a problem describes several things spreading or growing
outward at the same rate and asks for the time until the spread
finishes or the distance to the nearest source. The very next section,
01 Matrix, applies this exact same technique to a differently-phrased
problem.
\end{takeawaybox}

\section{Flood Fill}
\label{sec:graph-flood-fill}

\problemstatement{Given a 2D grid of pixel colors, a starting pixel
$(\textit{sr},\textit{sc})$, and a $\textit{newColor}$, change the
color of the starting pixel and every pixel \textbf{4-directionally
connected} to it through pixels of the \textbf{same original color} to
$\textit{newColor}$.}

\begin{patternbox}
\emph{``Change every connected pixel of the same color''} is a direct
single-source DFS (or BFS) on the implicit grid graph, where two
pixels are connected exactly when they are $4$-directional neighbors
\textbf{and} share the starting pixel's original color.
\end{patternbox}

\subsection*{Algorithm}

\begin{enumerate}
  \item Record $\textit{originalColor} = \textit{image}[\textit{sr}][\textit{sc}]$.
  \item If $\textit{originalColor} = \textit{newColor}$, return the
        image unchanged immediately (otherwise the DFS below would
        recurse infinitely, since every newly-colored pixel would
        still match the ``original color'' check).
  \item DFS from $(\textit{sr},\textit{sc})$: if the current pixel is
        in bounds and equals $\textit{originalColor}$, set it to
        $\textit{newColor}$, then recurse into all $4$ neighbors.
\end{enumerate}

\subsection*{C++ implementation}

\begin{lstlisting}
#include <bits/stdc++.h>
using namespace std;

void fill(vector<vector<int>>& image, int r, int c, int originalColor, int newColor) {
    int rows = image.size(), cols = image[0].size();
    if (r < 0 || r >= rows || c < 0 || c >= cols) return;
    if (image[r][c] != originalColor) return;

    image[r][c] = newColor;
    fill(image, r - 1, c, originalColor, newColor);
    fill(image, r + 1, c, originalColor, newColor);
    fill(image, r, c - 1, originalColor, newColor);
    fill(image, r, c + 1, originalColor, newColor);
}

vector<vector<int>> floodFill(vector<vector<int>>& image, int sr, int sc, int newColor) {
    int originalColor = image[sr][sc];
    if (originalColor != newColor) {
        fill(image, sr, sc, originalColor, newColor);
    }
    return image;
}

int main() {
    vector<vector<int>> image = {
        {1, 1, 1},
        {1, 1, 0},
        {1, 0, 1}
    };
    vector<vector<int>> result = floodFill(image, 1, 1, 2);
    for (auto& row : result) {
        for (int x : row) cout << x << " ";
        cout << endl;
    }
    // 2 2 2
    // 2 2 0
    // 2 0 1
    return 0;
}
\end{lstlisting}

\subsection*{Dry run}

\dryrun{Take the $3\times3$ image above, starting at $(1,1)$ (value
$1$), $\textit{newColor}=2$.}

$\textit{originalColor}=1 \neq 2$, proceed. Fill$(1,1)$: matches,
set to $2$; recurse to $(0,1)$: matches, set $2$; recurse further to
$(0,0)$ and $(0,2)$, both matching, set $2$. Back to $(1,1)$'s other
neighbors: $(2,1)=0$, no match, skip. $(1,0)=1$, matches, set $2$;
its neighbor $(2,0)=1$, matches, set $2$. $(1,2)=0$, no match, skip.
Final grid: top-left $2\times2$ block plus $(0,2)$ and $(2,0)$ all
become $2$, while the two $0$-pixels and the isolated $(2,2)=1$ pixel
(not $4$-connected to the start through same-colored pixels) remain
unchanged.

\begin{complexitybox}
\textbf{Time Complexity.} $O(\textit{rows} \times \textit{cols})$ --
each pixel is visited (and re-colored) at most once, since after being
set to $\textit{newColor}$ it no longer matches $\textit{originalColor}$
and further recursive calls on it return immediately.
\textbf{Space Complexity.} $O(\textit{rows} \times \textit{cols})$
recursion stack in the worst case (an entirely same-colored grid).
\end{complexitybox}

\begin{takeawaybox}
Flood Fill is the simplest possible instance of ``traverse an implicit
grid graph'' -- a single source, a single connectivity condition (same
original color), and no auxiliary bookkeeping beyond the grid itself
(the grid doubles as its own \textit{visited} array, since a re-colored
pixel naturally fails the match check on revisit).
\end{takeawaybox}

\section{01 Matrix}
\label{sec:graph-01-matrix}

\problemstatement{Given a binary matrix, return a matrix of the same
size where each cell holds the distance to its \textbf{nearest} cell
containing $0$ (distance measured in $4$-directional grid steps).}

\begin{patternbox}
\emph{``Distance to the nearest of several targets''} is
Section~\ref{sec:graph-rotten-oranges}'s multi-source BFS again, applied
to a distance-labeling problem instead of a time-to-completion problem:
seed the queue with every $0$-cell simultaneously (distance $0$), then
BFS outward, where each level reached corresponds to one more unit of
distance.
\end{patternbox}

\subsection*{Algorithm}

\begin{enumerate}
  \item Initialize a $\textit{dist}$ matrix to $0$ wherever the input
        is $0$, and push every such cell into the BFS queue.
  \item For every other cell, initialize $\textit{dist}$ to $-1$
        (unvisited marker) -- this doubles as the \textit{visited}
        check, exactly as a $-1$/sentinel value did in earlier
        chapters.
  \item Run BFS: for each cell popped, examine its $4$ neighbors; any
        neighbor still at $-1$ gets $\textit{dist}[\textit{neighbor}] =
        \textit{dist}[\textit{current}] + 1$ and is pushed.
  \item Return the filled $\textit{dist}$ matrix.
\end{enumerate}

\subsection*{C++ implementation}

\begin{lstlisting}
#include <bits/stdc++.h>
using namespace std;

vector<vector<int>> updateMatrix(vector<vector<int>>& mat) {
    int rows = mat.size(), cols = mat[0].size();
    vector<vector<int>> dist(rows, vector<int>(cols, -1));
    queue<pair<int,int>> q;

    for (int r = 0; r < rows; r++) {
        for (int c = 0; c < cols; c++) {
            if (mat[r][c] == 0) {
                dist[r][c] = 0;
                q.push({r, c});
            }
        }
    }

    int dr[] = {-1, 1, 0, 0};
    int dc[] = {0, 0, -1, 1};

    while (!q.empty()) {
        auto [r, c] = q.front();
        q.pop();

        for (int dir = 0; dir < 4; dir++) {
            int nr = r + dr[dir], nc = c + dc[dir];
            if (nr >= 0 && nr < rows && nc >= 0 && nc < cols &&
                dist[nr][nc] == -1) {
                dist[nr][nc] = dist[r][c] + 1;
                q.push({nr, nc});
            }
        }
    }
    return dist;
}

int main() {
    vector<vector<int>> mat = {
        {0, 0, 0},
        {0, 1, 0},
        {1, 1, 1}
    };
    vector<vector<int>> result = updateMatrix(mat);
    for (auto& row : result) {
        for (int x : row) cout << x << " ";
        cout << endl;
    }
    // 0 0 0
    // 0 1 0
    // 1 2 1
    return 0;
}
\end{lstlisting}

\subsection*{Dry run}

\dryrun{Take the $3\times3$ matrix above.}

Initial queue: every $0$-cell -- $(0,0),(0,1),(0,2),(1,0),(1,2)$, all
at $\textit{dist}=0$. Processing these, their unvisited neighbors get
$\textit{dist}=1$: $(1,1)$ (adjacent to $(0,1)$, $(1,0)$, and $(1,2)$,
whichever processes first), $(2,0)$ (adjacent to $(1,0)$), $(2,2)$
(adjacent to $(1,2)$). Next level: $(2,1)$, adjacent to both
$(1,1)$ (now $\textit{dist}=1$) and $(2,0)$/$(2,2)$ (also
$\textit{dist}=1$), gets $\textit{dist}=2$ (it is reached from a
$\textit{dist}=1$ cell first, since BFS processes level by level).
Final matrix matches the expected output, with the single $\textit{dist}=2$
cell at $(2,1)$ being the only cell two steps from the nearest $0$.

\begin{complexitybox}
\textbf{Time Complexity.} $O(\textit{rows} \times \textit{cols})$ --
each cell is pushed and processed at most once.
\textbf{Space Complexity.} $O(\textit{rows} \times \textit{cols})$ for
the \textit{dist} matrix and the queue.
\end{complexitybox}

\begin{takeawaybox}
Rotten Oranges and 01 Matrix are, structurally, the \emph{same}
algorithm wearing two different costumes: ``minutes until everything is
reached'' and ``distance to the nearest target'' are both just
\textbf{the BFS level number at which each cell is first reached},
read off in two different ways (the level of the very last cell
reached, versus the level of each individual cell respectively).
Recognizing multi-source BFS by its shape -- several simultaneous
starting points, uniform step cost -- matters more than the specific
wording of what the level number represents.
\end{takeawaybox}

\section{Number of Islands}
\label{sec:graph-number-of-islands}

\problemstatement{Given a 2D grid of \texttt{'1'} (land) and
\texttt{'0'} (water), count the number of \textbf{islands} -- maximal
groups of land cells connected $4$-directionally.}

\begin{patternbox}
This closing problem of the chapter is
Section~\ref{sec:graph-number-of-provinces}'s connected-component
counting loop, applied to an \textbf{implicit} grid graph instead of an
explicit adjacency matrix -- a land cell's neighbors are simply its
(up to) $4$ adjacent grid cells, with no adjacency list ever
constructed.
\end{patternbox}

\subsection*{Algorithm}

\begin{enumerate}
  \item Maintain a $\textit{visited}$ grid (or mutate the input grid
        directly, marking visited land as \texttt{'0'}) and an
        $\textit{islands}$ counter.
  \item For every cell: if it is land and unvisited, increment
        $\textit{islands}$, then DFS (or BFS) from it, marking every
        $4$-directionally connected land cell visited.
  \item Return $\textit{islands}$.
\end{enumerate}

\subsection*{C++ implementation}

\begin{lstlisting}
#include <bits/stdc++.h>
using namespace std;

void dfs(int r, int c, vector<vector<char>>& grid, vector<vector<bool>>& visited) {
    int rows = grid.size(), cols = grid[0].size();
    if (r < 0 || r >= rows || c < 0 || c >= cols) return;
    if (visited[r][c] || grid[r][c] == '0') return;

    visited[r][c] = true;
    dfs(r - 1, c, grid, visited);
    dfs(r + 1, c, grid, visited);
    dfs(r, c - 1, grid, visited);
    dfs(r, c + 1, grid, visited);
}

int numIslands(vector<vector<char>>& grid) {
    int rows = grid.size(), cols = grid[0].size();
    vector<vector<bool>> visited(rows, vector<bool>(cols, false));
    int islands = 0;

    for (int r = 0; r < rows; r++) {
        for (int c = 0; c < cols; c++) {
            if (grid[r][c] == '1' && !visited[r][c]) {
                islands++;
                dfs(r, c, grid, visited);
            }
        }
    }
    return islands;
}

int main() {
    vector<vector<char>> grid = {
        {'1', '1', '0', '0'},
        {'1', '1', '0', '0'},
        {'0', '0', '1', '0'},
        {'0', '0', '0', '1'}
    };
    cout << numIslands(grid) << endl; // 3
    return 0;
}
\end{lstlisting}

\subsection*{Dry run}

\dryrun{Take the $4\times4$ grid above.}

$(0,0)$: land, unvisited, $\textit{islands}=1$; DFS visits the
top-left $2\times2$ block of land. $(0,1),(1,0),(1,1)$: already
visited, skip. $(2,2)$: land, unvisited, $\textit{islands}=2$; DFS
visits just itself (no land neighbors). $(3,3)$: land, unvisited,
$\textit{islands}=3$; DFS visits just itself. Final answer: $3$
islands.

\begin{complexitybox}
\textbf{Time Complexity.} $O(\textit{rows} \times \textit{cols})$ --
every cell is visited at most once across all DFS calls combined.
\textbf{Space Complexity.} $O(\textit{rows} \times \textit{cols})$ for
the \textit{visited} grid and, in the worst case, the recursion stack
(an entirely land-filled grid).
\end{complexitybox}

\begin{takeawaybox}
This chapter's eight problems form a complete toolkit built from just
two primitives -- DFS and BFS -- applied to three settings (an explicit
adjacency list, an explicit adjacency matrix, and an implicit grid):
plain traversal, traversal with a hash map for building a copy,
component-counting via repeated traversal from unvisited starts, and
multi-source BFS for simultaneous spreading. Every later chapter in
this book's graph coverage layers a new constraint (direction, weight,
ordering) on top of exactly these same two traversal primitives.
\end{takeawaybox}

